\uuid{vOsS}
\exo7id{4181}
\titre{exo7 4181}
\auteur{quercia}
\organisation{exo7}
\datecreate{2010-03-11}
\isIndication{false}
\isCorrection{true}
\chapitre{Fonction de plusieurs variables}
\sousChapitre{Dérivée partielle}
\cpge{mpsi-mp2i.fonctions-deux-variables.a, pcsi.fonctions-deux-variables.a, ptsi.fonctions-deux-variables.a, mpsi-mp2i.fonctions-deux-variables.b, pcsi.fonctions-deux-variables.b, ptsi.fonctions-deux-variables.b, bcpst1.fonctions-deux-variables.b, bcpst1.fonctions-deux-variables.c, mp-mpi.topologie-espaces-normes.e, pc.espaces-vectoriels-normes.e, psi.espaces-vectoriels-normes.e, pt.fonctions-deux-trois-variables.b, pc.espaces-vectoriels-normes.f, psi.espaces-vectoriels-normes.f, pt.fonctions-deux-trois-variables.g}
\module{Analyse}
\niveau{L2}
\difficulte{}

\contenu{
\texte{
Soit $U$ un ouvert de~$\R^2$.
}
\begin{enumerate}
    \item \question{Donner un exemple de fonction $f : U \to \R$ ayant en tout point des
dérivées partielles premières, mais discontinue en au moins un point.}
\reponse{$xy/(x^2+y^2)$.}
    \item \question{Soit $f : U \to \R$ ayant en tout point des
dérivées partielles premières {\it bornées\/} sur~$U$. Montrer que~$f$
est continue.}
\reponse{$f$ est lipschitzienne pour $\|\ \|_\infty$.}
\end{enumerate}
}
